\section{Evaluation}

The empirical evaluation is intended to measure the cost of the primitive layer
and the end-to-end protocol separately. At the primitive level, the most
relevant metrics are the prover time, verifier time, and proof size of domain
lifting, replication, tensor sums, aggregation, and histogram construction. At
the application level, the evaluation reports the cost of batched inference and
end-to-end training proofs as functions of the number of trees, the maximum
depth, the feature dimension, the number of bins, and the batch size.

\subsection{Experimental Methodology}

The experiments compare the proposed protocol against the most relevant
baselines suggested by the literature and by the proof architecture itself.
Natural baselines include direct reuse of decision-tree-oriented lookup systems
\cite{cq++24}, data-parallel zkSNARK approaches to tree proving
\cite{sparrow24}, and recent training-oriented baselines such as
\cite{zkboost26}. Of particular interest is whether the histogram and
aggregation layer yields the intended improvement over proof strategies that
explicitly sort, compare, or RAM-simulate candidate splits inside a generic
backend.
